% Section 4.4, from lectures/L04/appendix/d_what_to_build.md.

\section{What to Build}
\label{c4:sec:build}

\subsection{The task}
\label{c4:sec:task}
Write \file{drivers/source/led_array.asm}, four subroutines that treat a run of LED structures as
one thing.

It is not large. What it is for is the first of the two skills this chapter is about, walking
memory at a stride that is not one. The second, knowing in advance that two things sharing 2048
bytes will not meet, is arithmetic you do on paper (\secref{c4:sec:stack}) and then check in the
simulator (\secref{c4:sec:exercises}).

Creating \file{drivers/source/led_array.asm} is all it takes to switch its tests on.

\subsection{The array}
\label{c4:sec:array}
An array of \code{count} LED structures, laid out end to end with no gaps, starting at an address
the caller chose. Entry $n$ therefore begins $n \times$ \code{LED_SIZE} bytes after the base, and
\code{LED_SIZE} is 7 (\file{led.inc}).

Nothing marks the end. The count is passed in, and it is the only thing that stops a walk.

\subsection{\code{led_array_init}}
\label{c4:sec:init}
\noindent\textbf{Arguments:} \code{r25:r24} the base address, \code{r22} the count, \code{r20} the
Arduino pin of the first LED.\\
\textbf{Returns:} \code{r24} is 0 on success, 1 if refused.

\medskip
Initialise \code{count} LEDs on \textbf{consecutive} Arduino pins, starting at the given one.
Entry 0 gets the first pin, entry 1 the next, and so on.

\noindent Write it like this:
\begin{enumerate}
  \item \textbf{Refuse an array that would run off the end of the pin numbering.} The last LED
    would be on pin \code{first + count - 1}, and there is no pin 14. Check this once, before
    initialising anything, rather than discovering it partway through and leaving half an array
    behind.
  \item \textbf{A count of zero succeeds and does nothing.} Test the count before the first
    iteration.
  \item \textbf{For each entry}, call \code{led_init} with the entry's address and the entry's pin
    number, then advance the pointer by \code{LED_SIZE} and the pin number by one.
\end{enumerate}

\subsubsection{Details that are pinned rather than up to you}
\textbf{\code{led_init} clobbers the argument registers, and it clobbers \code{Z}.} Your loop
variables cannot live in \code{r18} to \code{r27} or in \code{r30}/\code{r31} across the call, and
the pointer you are walking is in \code{r25:r24} precisely because that is where \code{led_init}
wants it. Pushing what you need before the call and popping it after is the straightforward answer,
and it costs four cycles per register per iteration; the alternative is to keep the loop state in
call-saved registers and pay to save those once, which is cheaper for a long array and more code.
Either is fine. Choosing without noticing there was a choice is not.

\textbf{Advance by \code{LED_SIZE}, taken from the include file.} Not by 7 written out, and
certainly not by 2 or 6, which are the sizes of the fields rather than of the structure. A stride of
6 puts each entry one byte on top of the last, and everything still looks like a structure.

\textbf{Do not call \code{led_init} for entries beyond the count.} The byte after the last structure
is an ordinary byte and will be initialised as happily as any other.

\textbf{Check yourself:} six LEDs from Arduino pin 8 gives entries whose pin fields are 0, 1, 2, 3,
4 and 5, all on port B, at addresses seven apart. Four LEDs from pin 6 straddle the ports: two on
port D at bits 6 and 7, then two on port B at bits 0 and 1.

\subsection{\code{led_array_all_on}, \code{led_array_all_off} and \code{led_array_toggle_at}}
\label{c4:sec:walks}
\textbf{\code{led_array_all_on}} and \textbf{\code{led_array_all_off}} take the base in
\code{r25:r24} and the count in \code{r22}, and call \code{led_on} or \code{led_off} for every
entry. They return nothing.

\textbf{\code{led_array_toggle_at}} takes the base in \code{r25:r24}, the count in \code{r22} and
an index in \code{r20}. It toggles that one entry and returns 0, or returns 1 without touching
anything if the index is not less than the count.

\subsubsection{Details that are pinned rather than up to you}
\textbf{The bound is \code{index < count}, and the interesting value is \code{index == count}.}
That is the first address past the end of the array; it looks like a structure, the driver will
accept it, and nothing else will notice. An implementation that checks \code{index <= count} passes
every test with a small index in it.

\textbf{Two of these three walk and one indexes.} Walking is a loop and an \code{adiw}; indexing is
a multiply, or a walk of \code{index} steps. Either is acceptable; if you use \code{mul}, clear
\code{r1} afterwards (\secref{c4:sec:stride}).

\textbf{Test the count at the top.} A count of zero means do nothing, for all three.

\subsection{What none of this does}
\label{c4:sec:notdo}
\textbf{It does not stop you.} The arithmetic in \secref{c4:sec:stack} produces a number and
nothing else. Nothing on the device checks your assembly against it, and a program whose stack
collides with its variables will assemble and run exactly as happily as one whose does not.

\textbf{It does not know about anything you did not count.} A subroutine that pushes two registers
adds two bytes that your worst case will not include unless you put them there. A bound is only as
good as the depth you hand it, which is the usual bargain and is why the exercises ask you to
measure as well.

\textbf{And the measurement does not know about the interrupt that has not arrived yet}, which is
why neither on its own is enough.
